% 妊神星（综述）
% license CCBYSA3
% type Wiki

本文根据 CC-BY-SA 协议转载翻译自维基百科\href{https://en.wikipedia.org/wiki/Haumea}{相关文章}。

帕洛玛天文台是一座位于美国加利福尼亚州圣迭戈县帕洛玛山的天文研究台站。该天文台由加州理工学院（Caltech）拥有并运营。观测时间由加州理工学院及其科研合作单位共享，这些单位包括喷气推进实验室（JPL）、耶鲁大学$^{[1]}$、以及中国科学院国家天文台$^{[2]}$。

该天文台运行多台望远镜，包括 200 英寸（5.1 米）的黑尔望远镜$^{[3]}$，48 英寸（1.2 米）的塞缪尔·奥申望远镜$^{[4]}$（用于 Zwicky 暂现天体设施 ZTF）$^{[5]}$，帕洛玛 60 英寸（1.5 米）望远镜$^{[6]}$，以及 30 厘米（12 英寸）的 Gattini-IR 望远镜$^{[7]}$。已退役的仪器包括帕洛玛试验台干涉仪，以及天文台最早启用的望远镜——建于 1936 年的 18 英寸（46 厘米）施密特相机。
\subsection{历史}
\begin{figure}[ht]
\centering
\includegraphics[width=6cm]{./figures/15c1f9cc6c3cbd5f.png}
\caption{帕洛玛山天文台出现在 1948 年的美国邮票上} \label{fig_Haumea_1}
\end{figure}
\subsubsection{Hale 对大型望远镜与帕洛玛天文台的愿景}
天文学家乔治·埃勒里·黑尔（George Ellery Hale），其愿景促成了帕洛玛天文台的建立，曾连续四次建造出当时世界上最大口径的望远镜$^{[8]}$。他在 1928 年发表了一篇文章，提出了后来成为 200 英寸帕洛玛反射望远镜的构想；这是一份向美国公众发出的邀请，旨在让他们了解大型望远镜如何帮助解答宇宙本质的基本问题。黑尔随后于 1928 年 4 月 16 日致信洛克菲勒基金会的国际教育委员会（后来并入一般教育委员会），在信中提出对此项目的资金申请。在信中，黑尔写道：

“在推动科学进步方面，没有任何方法能像开发全新的、更强大的仪器和研究方法那样富有成效。一台更大的望远镜不仅能够在空间穿透力与摄影分辨率上提供必要的提升，还可以让源自物理与化学最新基础性进展的理念与装置得以应用。”
\subsubsection{黑尔望远镜}
这台 200 英寸望远镜以天文学家兼望远镜制造者乔治·埃勒里·黑尔的名字命名。它由加州理工学院利用洛克菲勒基金会提供的 600 万美元资助建造，采用的是康宁玻璃厂在乔治·麦考利（George McCauley）指导下制造的 Pyrex 镜坯。J.A. Anderson 博士是最初的项目经理，其任命可追溯至 1930 年代早期$^{[9]}$。该望远镜（当时世界上最大）于 1949 年 1 月 26 日首次观测（first light），目标为 NGC 2261$^{[10]}$。美国天文学家埃德温·鲍威尔·哈勃（Edwin Powell Hubble）是第一位使用该望远镜的天文学家。

这台 200 英寸望远镜自 1949 年起一直保持世界最大望远镜的地位，直到 1975 年俄罗斯 BTA-6 望远镜首次观测。使用黑尔望远镜的天文学家在宇宙尺度上发现了类星体（后被归入所谓的活动星系核的一类）。他们研究了恒星族群的化学性质，从而理解了恒星核合成在宇宙中按观测丰度产生元素的起源，并已发现数千颗小行星。在纽约康宁市的康宁社区学院，存放着一台 1/10 比例的望远镜工程模型（康宁玻璃厂 — 即现今的康宁公司所在地），该模型曾用于发现至少一颗小行星——34419 Corning。
\subsubsection{建筑与设计}
\begin{figure}[ht]
\centering
\includegraphics[width=6cm]{./figures/6fd368c863632c16.png}
\caption{黑尔望远镜穹顶} \label{fig_Haumea_2}
\end{figure}
拉塞尔·W·波特（Russell W. Porter）设计了天文台建筑的装饰艺术（Art Deco）风格，包括 200 英寸黑尔望远镜的穹顶。波特还负责了黑尔望远镜与施密特相机的大部分技术设计工作，绘制了一系列横截面工程图。波特与许多工程师及加州理工学院委员会成员合作进行这些设计$^{[11]}$ $^{[12]}$ $^{[13]}$。
马克斯·梅森（Max Mason）负责施工，而西奥多·冯·卡门（Theodore von Karman）参与了工程设计。
\subsubsection{台长}
\begin{itemize}
\item Ira Sprague Bowen，1948–1964
\item Horace Welcome Babcock，1964–1978
\item Maarten Schmidt，1978–1980
\item Gerry Neugebauer，1980–1994
\item James Westphal，1994–1997
\item Wallace Leslie William Sargent，1997–2000
\item Richard Ellis，2000–2006
\item Shrinivas Kulkarni，2006–2018
\item Jonas Zmuidzinas，2018–2025
\item Mansi Kasliwal，2025–$^{[14]}$
\end{itemize}
\subsubsection{帕洛玛天文台与光污染}
南加州周边区域的大部分地区已采用遮光照明方式，以减少可能影响天文台的光污染$^{[15]}$。
\subsection{望远镜与仪器}
\begin{figure}[ht]
\centering
\includegraphics[width=6cm]{./figures/0287c69c9b783306.png}
\caption{黑尔望远镜穹顶} \label{fig_Haumea_3}
\end{figure}
\begin{figure}[ht]
\centering
\includegraphics[width=6cm]{./figures/5ac02f531ad92d9d.png}
\caption{黑尔望远镜的组成部件} \label{fig_Haumea_4}
\end{figure}
\begin{itemize}
\item 200 英寸（5.1 米）黑尔望远镜首次于 1928 年提出，并自 1949 年开始运行。它曾作为世界上最大的望远镜达 26 年之久$^{[3]}$（需要引用）。
\item 一台 60 英寸（1.5 米）反射望远镜位于 Oscar Mayer 大楼内，并实现完全机器人化运行。该望远镜于 1970 年投入使用，建造目的是为了提升帕洛玛天文学家对天空的观测覆盖率。其显著成果之一是发现了第一颗棕矮星$^{[6]}$。目前，该 60 英寸望远镜搭载 SED Machine 积分视场光谱仪，用于 ZTF 瞬变源的后续观测与分类$^{[16]}$ $^{[17]}$。
\item 48 英寸（1.2 米）塞缪尔·奥申望远镜的研制始于 1938 年，望远镜于 1948 年首次成像。它最初被称为 48 英寸施密特望远镜，并于 1986 年以塞缪尔·奥申命名$^{[4]}$。在众多重要成果中，奥申望远镜的观测促成了重要矮行星厄里斯（Eris）与赛德娜（Sedna）的发现$^{[18]}$。厄里斯的发现引发国际天文学界的讨论，最终导致冥王星于 2006 年被重新分类为矮行星。奥申望远镜目前实现完全机器人化运行，并搭载 5.7 亿像素的 ZTF 相机$^{[19]}$，作为 ZTF 项目的“发现引擎”。
\item 40 英寸（1.0 米）WINTER（宽视场红外瞬变探索者）1×1 度反射式机器人望远镜自 2021 年开始运行。其任务是针对红外（IR）天空进行视宁度受限的时间域巡天，特别关注由 LIGO 探测到的双中子星（BNS）并合残骸中 r 过程物质的识别。该仪器使用 Y、J 与短 H（Hs）波段观测，后者调整至 InGaAs 探测器的长波截止点，覆盖 0.9 至 1.7 微米的波长范围$^{[20]}$。
\end{itemize}
\subsubsection{已退役仪器}
\begin{itemize}
\item 一台 18 英寸（46 厘米）施密特相机于 1936 年成为帕洛玛天文台最早投入运行的望远镜。在 20 世纪 30 年代，弗里茨·兹威基（Fritz Zwicky）和沃尔特·巴德（Walter Baade）建议在帕洛玛增加巡天望远镜，于是开发了这台 18 英寸望远镜，用以展示施密特光学概念。兹威基利用这台 18 英寸望远镜在其他星系中发现了超过 100 颗超新星。1993 年，舒梅克-利维 9 号彗星（Shoemaker–Levy 9）也由这台仪器发现。该望远镜现已退役，并陈列在小型博物馆/游客中心中$^{[21]}$ $^{[22]}$。
\item 帕洛玛试验台干涉仪（Palomar Testbed Interferometer，PTI）是一台多望远镜组合仪器，能够对恒星的视直径和相对位置进行高角分辨率测量。PTI 测量了亮星的视直径及部分恒星的视形状，还测量了多星系统的视轨道。PTI 于 1995 年至 2008 年运行$^{[23]}$。
\item 帕洛玛行星搜索望远镜（Palomar Planet Search Telescope，PPST），又称 “Sleuth”，是一台 10 厘米（3.9 英寸）机器人望远镜，自 2003 年运行至 2008 年。它专用于通过凌日法搜索其他恒星周围的行星。该望远镜与洛厄尔天文台（Lowell Observatory）以及加那利群岛的望远镜联动，作为跨大西洋系外行星巡天（TrES）的一部分运行$^{[24]}$。
\end{itemize}
\subsection{研究}
\begin{figure}[ht]
\centering
\includegraphics[width=6cm]{./figures/fe35c967764026d3.png}
\caption{现已退役的 18 英寸施密特相机} \label{fig_Haumea_5}
\end{figure}
帕洛马天文台仍是一个活跃的科研设施，在每个晴天夜间运行多台望远镜，并支持一个规模庞大的国际天文学家群体从事广泛的研究领域。

海尔望远镜$^{[3]}$ 仍在开展活跃的科研观测，并在多个焦点配备多种光学及近红外光谱仪与成像相机。海尔望远镜还配备多级、高阶自适应光学系统，可在近红外波段提供衍射极限的成像能力。使用海尔望远镜取得的关键历史性科学成果包括：宇宙学尺度上的哈勃流测量、类星体作为活动星系核前身的发现、以及关于恒星族群与恒星核合成的研究。

奥申望远镜和 60 英寸望远镜以机器人方式运行，并共同支持主要的瞬变天文学项目——兹威基瞬变设施（ZTF）。

奥申望远镜建立的目的在于促进天文巡天，并用于许多重要的天文巡天项目，其中包括：
\subsubsection{POSS-I}
最初的帕洛马天文台巡天（POSS 或 POSS-I）由美国国家地理研究所赞助，完成于 1958 年。首批底片曝光于 1948 年 11 月，最后一批于 1958 年 4 月。该巡天使用奥申望远镜，以 14 英寸²（6 平方度）蓝敏（柯达 103a-O）和红敏（柯达 103a-E）照相底片进行观测。巡天覆盖从赤纬 +90°（天北极）到 −27° 的全天区，并覆盖全部赤经，灵敏度达到 22 等星（约为人眼可见极限的 100 万倍）。1957–1958 年，开展了 POSS 南部扩展，将巡天范围扩展至赤纬 −33°。最终的 POSS I 数据集包含 937 对底片。

数字化天空巡天（DSS）所生成的图像基于 POSS-I 巡天期间获取的胶片数据$^{[25]}$。

澳大利亚射电天文学家 J. B. Whiteoak 使用同一仪器将 POSS-I 数据向南扩展至赤纬 −42°。Whiteoak 的观测使用与对应北半球赤纬带相同的视场中心。与 POSS-I 不同，Whiteoak 扩展仅使用红敏（柯达 103a-E）照相底片。
\subsubsection{POSS-II}
“第二次帕洛马天文台巡天（Second Palomar Observatory Sky Survey）”重定向至此。关于 POSS I，参见 Palomar Sky Survey。

第二次帕洛马天文台巡天（POSS II，亦称 Second Palomar Sky Survey）于 20 世纪 80–90 年代实施，使用了更先进、更快速的胶片以及升级后的望远镜。奥申施密特望远镜升级为具备消色差改正镜，并提供自动导星功能。成像采用三个波长的照相底片：蓝光（IIIaJ，480 nm）、红光（IIIaF，650 nm）以及近红外（IVN，850 nm）。POSS II 的观测人员包括 C. Brewer、D. Griffiths、W. McKinley、J. Dave Mendenhall、K. Rykoski、Jeffrey L. Phinney，以及 Jean Mueller（她通过比较 POSS I 与 POSS II 底片发现了 100 余颗超新星）。在 POSS II 期间，Mueller 还发现了多颗彗星及小行星，明亮的威尔逊彗星（Comet Wilson 1986）即被当时还是研究生的 C. Wilson 在巡天早期发现$^{[26]}$。

在2 微米全天巡天（2MASS）完成之前，POSS II 一直是最广泛的宽视场天空巡天。完成后，斯隆数字巡天（SDSS）将在深度上超过 POSS I 与 POSS II，尽管 POSS 覆盖的天空面积约为 SDSS 的 2.5 倍。

POSS II 也以数字化形式存在（即对照相底片进行扫描），作为数字化天空巡天（DSS）的一部分$^{[27]}$。
\subsubsection{QUEST}
长期的 POSS 项目之后，开展了帕洛马类星体赤道巡天团队（Palomar Quasar Equatorial Survey Team, QUEST）的变源巡天$^{[28]}$。这一巡天成果被用于多个项目，包括近地小行星跟踪（NEAT）。另一项利用 QUEST 数据的计划于 2003 年 11 月 14 日发现了 90377 Sedna，并发现了约 40 个柯伊伯带天体。其他共享该照相机的项目包括：Shrinivas Kulkarni 的伽马射线爆搜索（利用自动望远镜能在爆发被探测到后立即反应，并对衰减期间连续拍摄）、Richard Ellis 的超新星搜索（用于检验宇宙膨胀是否在加速）、以及 S. George Djorgovski 的类星体搜索。

用于帕洛马 QUEST 巡天的照相机由 112 个电荷耦合器件（CCD）组成马赛克阵列，覆盖整个施密特望远镜的视场（4° × 4°）。在建造时，它是当时天文照相机中规模最大的 CCD 马赛克仪器。该仪器曾被用于制作《The Big Picture》，这是迄今制作的最大天文照片$^{[29]}$。《The Big Picture》目前陈列于格里菲斯天文台。
\subsubsection{当前研究}
目前使用 200 英寸海尔望远镜的研究项目覆盖可观测宇宙的广泛领域，包括近地小行星、外太阳系行星、柯伊伯带天体、恒星形成、系外行星$^{[30]}$、黑洞与 X 射线双星、超新星与其他瞬变源后续观测，以及类星体/活动星系核研究$^{[31]}$。

48 英寸塞缪尔·奥申施密特望远镜以机器人方式运行，并支持新的瞬变天文巡天——兹威基瞬变设施（ZTF）$^{[5]}$。

60 英寸望远镜以机器人方式运行，并通过指定用途的 光谱能量分布机（SEDM）$^{[32]}$ 积分视场光谱仪，提供快速、低色散的光学光谱，用于瞬变初步分类，从而支持 ZTF。
\subsection{参观与公众参与}
\begin{figure}[ht]
\centering
\includegraphics[width=6cm]{./figures/b2f47b475a9262d1.png}
\caption{“帕洛马天文台的格林威游客中心，内设礼品店”} \label{fig_Haumea_6}
\end{figure}
帕洛马天文台是一座活跃的科研设施，但部分区域在白天向公众开放。访客可在每天上午 9 点至下午 3 点自行参观 200 英寸望远镜。天文台全年每周 7 天开放，除 12 月 24 日、25 日及恶劣天气期间外。4 月至 10 月的周六和周日提供 200 英寸哈雷望远镜穹顶和观测区的导览参观。公众还可通过社区支持组织“帕洛马天文台讲解员”（Palomar Observatory Docents）参与幕后导览$^{[33]}$。

帕洛马天文台设有一座现场博物馆——“格林威游客中心”（Greenway Visitor Center）$^{[22]}$，内设与天文台及天文学相关的展览、礼品店$^{[34]}$，并定期举办公众活动$^{[35]}$。

对于无法前往天文台的人士，帕洛马提供了一套内容丰富的虚拟参观系统，可虚拟访问园区内所有主要科研望远镜和格林威游客中心，并包含大量嵌入式多媒体资料以提供更多背景信息$^{[36]}$。此外，天文台还积极维护官方网站$^{[37]}$和 YouTube 频道$^{[38]}$以促进公众参与。

天文台位于加利福尼亚州圣迭戈县北部 76 号州道旁，距圣迭戈市中心车程约两小时，距洛杉矶市中心（UCLA、LAX 机场）车程约三小时$^{[39]}$。入住附近帕洛马露营地的游客还可徒步 2.2 英里（3.5 公里）沿“天文台步道”抵达帕洛马天文台$^{[40]}$。值得一提的是，《星球大战》原版的音效设计师本·伯特（Ben Burt）曾在帕洛马天文台录制多种声音，包括穹顶电机及百叶窗声，用作“死星”背景音效。
\subsection{气候}
帕洛马属炎热夏季地中海气候（柯本气候分类 Csa）。
\begin{figure}[ht]
\centering
\includegraphics[width=6cm]{./figures/bdb5238781d002ad.png}
\caption{} \label{fig_Haumea_7}
\end{figure}